% Volume I, appendice «Elementi di statica grafica»: figura fig:nodi_sgr.
% Richiamata con \FiguraSGR{figura-sgr-nodi}; didascalia e label nell'appendice.
% Generata con gen_sgr.py (coordinate numeriche; stili di Preambolo/tikz-pol e tikz-sgr).
%
% Forze (unita' di disegno del poligono delle forze): f_1 = (-0.6, -1.2), f_2 = (0.1, -1.0), f_3 = (1.4, -1.0);
% vertici 0 = (0, 0), 1 = (-0.600,-1.200), 2 = (-0.500,-2.200), 3 = (0.900,-3.200); polo P = (2.200,-1.500).
% Nodi: A_1 = (1.864,-1.271), A_2 = (3.848,-1.484), A_3 = (6.672,-0.752); lato b_i parallelo al raggio iP.
% Forze nei lati N_i = |iP|: 2.663, 2.816, 2.789, 2.140.
% Nel piano di posizione le forze sono disegnate in scala 0.5 rispetto al poligono delle forze.
% (a) poligono funicolare materializzato: forze f_i applicate nei nodi (coda nel nodo), forze di
%     estremita' f_A = P0 in A_0 e f_B = 3P in A_4; campi: P dalla parte opposta alle forze, 0 ... 3
%     dalla parte delle forze, separati dalle rette d'azione (la forza f_i separa i campi i-1 e i,
%     il lato b_i separa i campi i e P);
% (b) poligono delle forze chiuso P 0 1 2 3 P; i triangoli (i-1, i, P), a grigi alterni, sono i
%     poligoni di equilibrio dei nodi A_i; sui raggi interni 1P e 2P le due forze opposte che il
%     lato esercita sui due nodi di estremita'.
\begin{tikzpicture}[x=1cm,y=1cm,font=\footnotesize]
\begin{scope}
\draw[polretta] (1.864,-1.271) -- (1.341,-2.318);
\draw[polforza] (1.864,-1.271) -- (1.564,-1.871);
\node[anchor=east,inner sep=3pt] at (1.654,-1.691) {$\fb_1$};
\draw[polretta] (3.848,-1.484) -- (3.948,-2.481);
\draw[polforza] (3.848,-1.484) -- (3.898,-1.984);
\node[anchor=west,inner sep=3pt] at (3.883,-1.834) {$\fb_2$};
\draw[polretta] (6.672,-0.752) -- (7.779,-1.542);
\draw[polforza] (6.672,-0.752) -- (7.372,-1.252);
\node[inner sep=0pt] at (6.883,-1.197) {$\fb_3$};
\draw[polestremo] (0.666,-0.454) -- (-0.434,0.296) node[pos=0.7,anchor=south west,inner sep=2pt] {$\fb_A$};
\draw[polestremo] (7.371,0.162) -- (8.021,1.012) node[pos=0.7,anchor=north west,inner sep=2pt] {$\fb_B$};
\draw[pollato] (0.666,-0.454) -- (1.864,-1.271) -- (3.848,-1.484) -- (6.672,-0.752) -- (7.371,0.162);
\pic at (0.666,-0.454) {polpunto};
\pic at (1.864,-1.271) {polpunto};
\pic at (3.848,-1.484) {polpunto};
\pic at (6.672,-0.752) {polpunto};
\pic at (7.371,0.162) {polpunto};
\node[anchor=north east,inner sep=2.5pt] at (0.686,-0.494) {$A_0$};
\node[anchor=south west,inner sep=2.5pt] at (1.814,-1.221) {$A_1$};
\node[anchor=south,inner sep=4pt] at (3.848,-1.484) {$A_2$};
\node[anchor=south east,inner sep=2.5pt] at (6.672,-0.712) {$A_3$};
\node[anchor=north west,inner sep=2.5pt] at (7.351,0.122) {$A_4$};
\node[inner sep=0pt] at (1.389,-0.681) {$b_0$};
\node[inner sep=0pt] at (2.979,-1.169) {$b_1$};
\node[inner sep=0pt] at (5.205,-0.905) {$b_2$};
\node[inner sep=0pt] at (6.812,-0.207) {$b_3$};
\node[polcampo] at (0.783,-1.442) {$0$};
\node[polcampo] at (2.771,-2.173) {$1$};
\node[polcampo] at (5.417,-1.955) {$2$};
\node[polcampo] at (7.652,-0.705) {$3$};
\node[polcampo] at (4.368,0.012) {$P$};
\node at (4.000,-3.034) {(a)};
\end{scope}
\begin{scope}[shift={(9.9,1.2)}]
\fill[black!9] (0.000,0.000) -- (-0.600,-1.200) -- (2.200,-1.500) -- cycle;
\fill[black!20] (-0.600,-1.200) -- (-0.500,-2.200) -- (2.200,-1.500) -- cycle;
\fill[black!9] (-0.500,-2.200) -- (0.900,-3.200) -- (2.200,-1.500) -- cycle;
\draw[polraggio] (-0.600,-1.200) -- (2.200,-1.500);
\draw[polraggio] (-0.500,-2.200) -- (2.200,-1.500);
\draw[polestremo] (2.200,-1.500) -- (0.000,0.000) node[pos=0.45,anchor=south west,inner sep=2pt] {$\fb_A$};
\draw[polestremo] (0.900,-3.200) -- (2.200,-1.500) node[pos=0.55,anchor=north west,inner sep=2pt] {$\fb_B$};
\draw[polforza] (0.000,0.000) -- (-0.600,-1.200) node[midway,anchor=east,inner sep=2.5pt] {$\fb_1$};
\draw[polforza] (-0.600,-1.200) -- (-0.500,-2.200) node[midway,anchor=east,inner sep=2.5pt] {$\fb_2$};
\draw[polforza] (-0.500,-2.200) -- (0.900,-3.200) node[midway,anchor=north east,inner sep=2.5pt] {$\fb_3$};
\draw[polverso] (0.377,-1.214) -- (0.794,-1.259);
\draw[polverso] (0.775,-1.438) -- (0.358,-1.393);
\draw[polverso] (0.408,-1.872) -- (0.815,-1.766);
\draw[polverso] (0.860,-1.940) -- (0.453,-2.046);
\node[inner sep=0pt,font=\scriptsize] at (0.217,-0.786) {$A_1$};
\node[inner sep=0pt,font=\scriptsize] at (0.018,-1.659) {$A_2$};
\node[inner sep=0pt,font=\scriptsize] at (0.613,-2.452) {$A_3$};
\pic at (0.000,0.000) {polpunto};
\pic at (-0.600,-1.200) {polpunto};
\pic at (-0.500,-2.200) {polpunto};
\pic at (0.900,-3.200) {polpunto};
\pic at (2.200,-1.500) {polpolo};
\node[anchor=south,inner sep=3.5pt] at (0.000,0.000) {$0$};
\node[anchor=east,inner sep=3.5pt] at (-0.600,-1.200) {$1$};
\node[anchor=east,inner sep=3.5pt] at (-0.500,-2.200) {$2$};
\node[anchor=north,inner sep=3.5pt] at (0.900,-3.200) {$3$};
\node[anchor=west,inner sep=3.5pt] at (2.200,-1.500) {$P$};
\node at (0.900,-4.234) {(b)};
\end{scope}
\end{tikzpicture}%
